% A.65 candidate; notes only, pending independent proof/source/code audits.
\subsection{Every intermediate cover of an odd Heisenberg cover}

\begin{statement}{A.65}{Open}{Candidate: all intermediate exponent-prime class-two covers}
Let \(p\) be any odd prime, and let
\(Y_p\to U\) be the actual degree-\(p^3\)
Heisenberg cover of A.61--A.62. Every geometrically
connected intermediate cover of \(Y_p\to U\)
satisfies both assertions of Task A.5 at every
coefficient prime, for every cusp packet and every
permitted open derived subgroup.

More intrinsically, this holds for every finite cover
of \(U\) whose geometric Galois closure has exponent
\(p\) and nilpotency class at most two. The new
nonabelian intermediate type has degree \(p^2\),
proper genus \(1+p(p-1)/2\), and \(p\) cusps.
Its closure is the full Heisenberg group. This does
not settle all nilpotent covers or the universal
Task A.5.
\end{statement}

\begin{proof}[Candidate proof, pending independent audits]
\paragraph{The actual subgroup alternatives.}
Write \(C=C_p\), \(q:C\to E\),
\(G=H_p\), and \(Z=Z(G)\). The actual map
to \(U\) is \([p]q\), and
\[
  1\longrightarrow Z\longrightarrow G
       \xrightarrow{\;\tau\;} E[p]\longrightarrow1,
  \qquad C/Z=E.                                      \tag{A.65.1}
\]
The commutator on \(E[p]\simeq\mathbf F_p^2\)
is a nondegenerate alternating form with values
in \(Z\simeq\mathbf F_p\). If a subgroup
\(S\le G\) contains \(Z\), its proper quotient
is \(E/(S/Z)\). The corresponding affine cover
has genus one and is covered by A.40.

If \(S\cap Z=1\), then \([S,S]=1\) and
the image \(\tau(S)\) is isotropic. In dimension
two an isotropic subspace has dimension at most
one. Thus \(S=1\), the case of A.61--A.62, or
\(S\) is a noncentral subgroup of order \(p\).
We treat the latter without replacing its quotient
by a pro-prime inclusion.

Every geometric point stabilizer for the actual
action of \(G\) on \(C\) is either trivial or
\(Z\). Indeed an element with nonzero translation
\(\tau(g)\) cannot fix a point over the elliptic
curve, and the cyclic degree-\(p\) map \(q\)
is totally ramified exactly above \(E[p]\).
Consequently \(S\) acts freely on the entire
generic proper curve. For
\[
       \pi:C\longrightarrow D=C/S
\]
we have an unramified degree-\(p\) map and
\[
 g(D)=1+\frac{p(p-1)}2,\qquad
 \#\{\text{cusps of }D\}=p.                         \tag{A.65.2}
\]
The genus follows from Riemann--Hurwitz and
\(g(C)=1+p^2(p-1)/2\); the cusp count follows
from the free \(S\)-action on the \(p^2\)
cusps. In particular there are at least three
distinct cusps.

\paragraph{A genuine elliptic quotient square.}
Put \(T=\tau(S)\), let
\(\phi:E\to A=E/T\) be the degree-\(p\)
isogeny, and let \(\widehat\phi:A\to E\)
be its dual. Actual equivariance of \(q\)
gives a degree-\(p\) map \(q_D:D\to A\)
with
\[
       q_D\pi=\phi q,\qquad
       \widehat\phi q_D\pi=[p]q.                    \tag{A.65.3}
\]
Thus the intermediate affine map to \(U\) is
\(\widehat\phi q_D\), of degree \(p^2\).
Its removed set is exactly
\(q_D^{-1}(\ker\widehat\phi)\). There is one
point above each of the \(p\) points of
\(\ker\widehat\phi\), with ramification
index \(p\). All these curves, groups, maps,
and marked points are defined over a suitable
finite number field. We may make subsequent
finite extensions without changing the geometric
subgroup or the selected packet.

\paragraph{Ordinary primes: choose a place with nonzero reduced translation.}
Suppose \(p\equiv1\pmod4\), and factor
\((p)=(\varpi)(\bar\varpi)\) in
\(\mathbf Z[i]\). The two lines
\(E[\varpi]\) and \(E[\bar\varpi]\)
are distinct in \(E[p]\). At a place above
\(\varpi\), the connected kernel of reduction
on \(E[p]\) is \(E[\varpi]\), while at a
place above \(\bar\varpi\) it is
\(E[\bar\varpi]\). To recall the exact
reason, the Gaussian endomorphisms extend to
the good elliptic model. The derivative of
\([\varpi]\) on an invariant differential is
\(\varpi\). Its degree is \(p\), so its
special map is purely inseparable at the first
place and etale at the second. The other factor
has the reverse behavior. These are the formal
CM identifications used in A.40, with
\cite[Theorem III.6.2 and Theorem V.3.1]{fghSilvermanAEC}.

Since \(T\) is one line, it cannot equal both
connected lines. Choose a place where it does
not equal the connected line. Then every
nonzero element of \(T\) has nonzero reduction.
This remains true after further finite extension.

Use the actual A.62 semistable model, with its
map to the good elliptic model and with all
original cusps separated. The \(G\)-action
extends to it; one may take the stable marked
model and then an equivariant resolution. For
each \(1\ne s\in S\), equivariance gives
\[
       \bar q(sP)=\bar q(P)+\overline{\tau(s)}
                         \ne\bar q(P)                \tag{A.65.4}
\]
at every geometric special point. Thus \(S\)
has no special stabilizer, including at nodes.
There is also no generic stabilizer. The finite
constant-group quotient is therefore an etale
\(S\)-torsor, even though \(|S|=p\): the
constant group scheme \(\mathbf Z/p\) is etale.
The quotient is a semistable proper model of
\(D\), by etale descent.

A.62 has \(p\) rational marked clusters above
the \(p\) distinct reduced points of \(E[p]\),
and each cluster contains \(p\) cusps. The
reduced translations of \(T\) permute these
clusters freely and transitively. Choose a lift
of any selected cusp of \(D\), and choose a
different cusp in its same rational cluster.
The two lifts are not \(S\)-equivalent, because
every nonidentity element changes the cluster.
Their rational joining path maps to a rational
joining path between distinct quotient cusps:
the quotient identifies the \(p\) copies of
this cluster and is an isomorphism on each one.
Thus A.42 gives the full fixed-inertia conclusion
at every coefficient prime.

\paragraph{Inert primes: compact type descends and no new positive component is hidden.}
Now suppose \(p\equiv3\pmod4\). The elliptic
curve \(E\) has good supersingular reduction.
After finite extension so does its isogenous
curve \(A\). The maps \(\phi\) and
\(\widehat\phi\) extend to the good group
models and retain degree \(p\): the degree
of the pullback of a relative degree-one
line bundle is constant. Their special maps
are purely inseparable, since a supersingular
elliptic curve has no nonzero geometric
\(p\)-torsion. Every point of
\(\ker\widehat\phi\) specializes to \(O\).

The norm on Jacobians is surjective, and
\(\pi_*\pi^*=[p]\). The potentially good
reduction of \(J_C\), supplied by A.61,
therefore implies potential good reduction
of \(J_D\), by
\cite[Theorem 1 and Corollary 2, p.493]{fghSerreTate1968}.
After choosing semistable reduction of \(D\),
its dual graph is consequently a tree. More
explicitly, the special identity component
of its Neron model is the special \(\operatorname{Pic}^0\),
and its torus rank is the graph cycle rank;
good Jacobian reduction forces this rank to
vanish. The exact comparison and rank formula
are recalled in \cite[Remark 10.4.24 and
Lemma 7.5.18(b)]{fghLiu}.

We justify the required actual model maps.
Start with regular semistable models and
resolve any rational map to a projective
target by point blowups. The relevant precise
result is \cite[Theorem 9.2.7, pp.396--397;
Proposition 9.2.5]{fghLiu}; every added
component is a projective line. In particular
the rational map from a regular model of
\(D\) to the good elliptic model \(A\)
extends without adding a positive component,
and in fact descends to the original model.
Here is the descent argument. At the last
blowup its exceptional projective line maps
constantly to the elliptic special fibre.
The image of the graph of the resolved map
is proper and quasi-finite over the original
normal surface: above the centre its fibre
is one point, and elsewhere it is one point.
It is therefore finite birational and hence
an isomorphism. Descending successively proves
the claimed extension of \(q_D\). This
argument avoids applying the Neron mapping
property to a source that is not smooth
over the valuation ring.

Next resolve the actual \(\pi:C\to D\)
by point blowups of the regular source model.
The resulting proper map is surjective: its
closed image contains the entire generic
fibre of the integral target model. Every
positive-genus component of the target
special fibre is therefore dominated by
one of the original positive-genus source
components. Indeed the added components are
rational, and a projective line cannot
dominate a positive-genus normalization,
by Luroth's theorem. This proves the
comparison for every positive target
component, including ones whose appearance
might otherwise be obscured by a wild quotient.
We have made no assertion that wild quotient
commutes with reduction.

A.61 has one elliptic main component whose
map to \(E\) is radicial. All other
positive-genus source components map
constantly to finite zeros of its function
\(R\), none of which is \(O\). Since the
special \(\phi\) is radicial, their images
in \(A\) are still different from \(O\).
The composite from the source main to
\(A\) is radicial and nonconstant. In
the resolved diagram (A.65.3) it factors
through precisely one target component,
which we call \(D_0\). Both factors are
radicial: their function fields are
intermediate fields of a purely inseparable
extension. In particular \(D_0\) has
exactly one point above \(O\).

No other positive target component maps
nonconstantly to \(A\), since the unique
source main cannot dominate two components
and all source tails have constant images.
Every other positive target component
maps constantly to a point different from
\(O\), by the domination argument. Also
\(D_0\) has genus one: it is dominated
by an elliptic curve and maps nonconstantly
to an elliptic curve.

\paragraph{All quotient cusps have rational joining paths.}
Every cusp of \(D\) specializes above
\(O\in A\). It cannot lie on any of the
other positive-genus components. A rational
component carrying such a mark has constant
map \(O\). Its path in the special-fibre
tree to \(D_0\) cannot pass through a
positive component with a different
constant image: equality of images at
each intervening node would be violated.
All such rational parts attach to the
unique point of \(D_0\) above \(O\).
At a nodal point only one other branch can
meet this component, so these parts form
one connected rational subtree.

If this subtree is absent before marking,
the \(p\ge3\) cusps all specialize to
the same smooth point of \(D_0\), and
their separation inserts a rational tree.
If it is present, its attachment on
\(D_0\) is a node; disjoint smooth marked
sections lie on the rational subtree or
on rational components inserted above
its nodes. In either case every selected
cusp has a distinct companion with a
fully rational joining path. A.41--A.42
give the desired full integral fixed
group for every coefficient prime and
every packet. Finite field restriction
inside a prescribed open derived subgroup
is exactly the one in A.41.

Every intermediate cover has nilpotent
geometric Galois closure, so A.59 gives
the normalizer assertion at all coefficient
primes. We have used no transfer of fixed
subgroups through a surjective pro-prime
quotient.

\paragraph{Intrinsic closure formulation.}
The geometric fundamental group of \(U\)
is free profinite on two generators. Its
maximal class-two exponent-\(p\) quotient
has the presentation
\[
 \langle x,y,z\mid z=[x,y],\ z\text{ central},
                      \ x^p=y^p=z^p=1\rangle.         \tag{A.65.5}
\]
Every word has the form \(x^a y^b z^c\),
with exponents modulo \(p\); the upper
unitriangular matrix model realizes
\(p^3\) distinct words. For odd \(p\),
that matrix group has exponent \(p\),
so (A.65.5) is exactly \(H_p\).
Every nonabelian quotient still has a
two-dimensional abelianization and
nontrivial one-dimensional commutator,
hence has order \(p^3\) and is \(H_p\).
The kernel of any surjection to \(H_p\)
is consequently the same characteristic
kernel of this universal quotient.

The actual A.61--A.62 cover realizes that
quotient. Thus every cover with geometric
closure as in the statement is geometrically
an intermediate of this one, or has abelian
closure and genus one. Any isomorphism and
finite descent data are defined after a
finite number-field extension. The preceding
all-packet, open-derived conclusions pass
to the original arithmetic setting exactly
as in A.41. This proves the stated scope
without a cofinality assertion.

The executed certificate
\path{notes/fgh_A65_certificate.py}
checks the exact finite Heisenberg subgroup
classification, cusp orbits, and the two
ordinary connected-line choices. It is
support for the uniform proof, not a
substitute for the model comparison.
\end{proof}
